% Source printed pp. 105-108. Subsections 4.1.1-4.1.4.
\section{代数空间与代数叠的定义}\label{src:4.1}
\begin{epigraph}
即使最抽象的定义, 一旦你熟悉了它, 也就不再抽象了. 它就像一座美丽的山, 你能看得十分清楚, 因为空气非常澄澈, 光线让你看见所有细节.
\epigraphauthor{Claire Voisin}
\end{epigraph}
什么是代数空间, Deligne--Mumford 叠和代数叠? 给出定义之后, 我们将验证商叠 $[U/G]$, 曲线的模叠 $\mathcal M_g$ 以及向量丛的模叠 $\cBun_{r,d}(C)$ 的代数性.

\subsection{代数空间}\label{src:4.1.1-subsection}
\begin{definition}[可由概形表出的态射]\label{src:4.1.1}
$\Sch$ 上的预叠 (或预层) 态射 $\mathcal X\to\mathcal Y$ 称为可由概形表出 (有时也称为概形态射 (schematic)), 如果对从概形出发的每个态射 $T\to\mathcal Y$, 纤维积 $\mathcal X\times_{\mathcal Y}T$ 都是概形.

若 $P$ 是在基变换下稳定的概形态射性质 (如满射性或平展性), 则称可由概形表出的预叠态射 $\mathcal X\to\mathcal Y$ 具有性质 $P$, 如果对从概形出发的每个态射 $T\to\mathcal Y$, 概形态射 $\mathcal X\times_{\mathcal Y}T\to T$ 都具有性质 $P$.
\end{definition}
\begin{definition}[代数空间]\label{src:4.1.2}
\newterm{代数空间}{algebraic space}是 $\Sch_{\text{ét}}$ 上的层 $X$, 使得存在概形 $U$ 以及可由概形表出的满射平展态射 $U\to X$.

映射 $U\to X$ 称为\newterm{平展呈示}{étale presentation}. 代数空间的态射按定义就是层的态射. 每个概形都是代数空间.
\end{definition}

\subsection{Deligne--Mumford 叠}\label{src:4.1.2-subsection}
\begin{definition}[可表态射及其性质]\label{src:4.1.3}
$\Sch$ 上的预叠 (或预层) 态射 $\mathcal X\to\mathcal Y$ 称为\newterm{可表的}{representable}, 如果对从概形 $T$ 出发的每个态射 $T\to\mathcal Y$, 纤维积 $\mathcal X\times_{\mathcal Y}T$ 都是代数空间.

设 $P$ 是在基变换下稳定且在源上平展局部的概形态射性质 (即若 $X'\to X$ 是满射平展态射, 则概形态射 $X\to Y$ 具有 $P$ 当且仅当 $X'\to X\to Y$ 具有 $P$); 例如满射性, 平展性或光滑性. 称可表的预叠态射 $\mathcal X\to\mathcal Y$ 具有性质 $P$, 如果对从概形出发的每个态射 $T\to\mathcal Y$, 以及由概形给出的每个平展呈示 $U\to\mathcal X\times_{\mathcal Y}T$, 复合 $U\to\mathcal X\times_{\mathcal Y}T\to T$ 都具有性质 $P$.
\end{definition}
\begin{definition}[Deligne--Mumford 叠]\label{src:4.1.4}
\newterm{Deligne--Mumford 叠}{Deligne--Mumford stack}是 $\Sch_{\text{ét}}$ 上的叠 $\mathcal X$, 使得存在概形 $U$ 及满射, 平展且可表的态射 $U\to\mathcal X$.

态射 $U\to\mathcal X$ 称为平展呈示. Deligne--Mumford 叠的态射按定义就是叠的态射. 由例~\bookref{3.4.7}, 每个代数空间都是 Deligne--Mumford 叠. 后面将证明, 若一个 Deligne--Mumford 叠 (甚至代数叠) 是层, 则它是代数空间 (定理~\bookref{4.6.6}~和~\bookref{5.5.10}).
\end{definition}
\begin{remark}\label{src:4.1.5}
代数空间与 Deligne--Mumford 叠的本质区别是前者为层, 后者为叠; 此外还有一个技术上的区别: 代数空间的平展呈示可由概形表出, 而 Deligne--Mumford 叠的平展呈示只要求可表. 若 Deligne--Mumford 叠的对角线分离且拟紧, 则对角线可由概形表出, 并且每个呈示 $U\to\mathcal X$ 都可由概形表出 (推论~\bookref{5.5.8}). 要求对角线分离且拟紧是十分温和的条件, 几乎所有值得研究的模问题都满足它. 但应注意, 确实存在对角线不拟紧, 不分离或不可由概形表出的 Deligne--Mumford 叠, 例如 $BG$, 其中 $G$ 是不拟紧, 不分离或不是概形的平展群代数空间; 见例~\bookref{4.9.36}~和~\bookref{4.9.38}.
\end{remark}

\subsection{代数叠}\label{src:4.1.3-subsection}
\begin{definition}[代数叠]\label{src:4.1.6}
\newterm{代数叠}{algebraic stack}是 $\Sch_{\text{ét}}$ 上的叠 $\mathcal X$, 使得存在概形 $U$ 及满射, 光滑且可表的态射 $U\to\mathcal X$.

态射 $U\to\mathcal X$ 称为\newterm{光滑呈示}{smooth presentation}. 代数叠的态射按定义就是预叠的态射. 每个概形, 代数空间或 Deligne--Mumford 叠都是代数叠.
\end{definition}
\begin{caution}\label{src:4.1.7}
文献中的多数作者会对对角线附加可表性条件. 然而, 如定理~\bookref{4.2.1}~所示, 光滑呈示的存在性蕴含对角线的可表性: 代数空间的对角线可由概形表出, 而代数叠的对角线可表.
\end{caution}
\begin{definition}[开子叠与闭子叠]\label{src:4.1.8}
$\Sch_{\text{ét}}$ 上的叠 $\mathcal X$ 的子叠 $\mathcal T\subseteq\mathcal X$ 称为\newterm{开子叠}{open substack} (分别为\newterm{闭子叠}{closed substack}), 如果包含 $\mathcal T\to\mathcal X$ 可由概形表出且为开浸入 (分别为闭浸入).
\end{definition}
\begin{exercise}[纤维积]\label{src:4.1.9}
证明代数空间, Deligne--Mumford 叠和代数叠的纤维积均存在.
\end{exercise}

\subsection{商叠的代数性}\label{src:4.1.4-subsection}
\begin{theorem}[商叠的代数性]\label{src:4.1.10}
若 $G\to S$ 是作用于代数空间 $U\to S$ 的光滑仿射群概形, 则商叠 $[U/G]$ 是 $S$ 上的代数叠, 且 $U\to[U/G]$ 是主 $G$-丛, 特别地, 它满射, 光滑且仿射. 特别地, 分类叠 $BG=[S/G]$ 是代数的.
\end{theorem}
\begin{remark}\label{src:4.1.11}
细心的读者可能已经注意到, 我们仅在 $U$ 为概形时定义了 $[U/G]$. 推广这一定义并不困难. 光滑仿射群概形 $G\to S$ 在代数空间 $U\to S$ 上的作用, 是满足与定义~\bookref{B.1.10}~相同公理的态射 $\sigma:G\times_S U\to U$.
我们将商叠 $[U/G]$ 定义为预叠 $[U/G]^{\mathrm{pre}}$ 的叠化, 后者在 $S$-概形 $T$ 上的纤维范畴为商群胚 $[U(T)/G(T)]$, 如定义~\bookref{3.4.15}~所述.
$[U/G]$ 在 $S$-概形 $T$ 上的对象是主 $G$-丛 $P\to T$ 及 $G$-等变态射 $P\to U$. 由于到代数空间的态射在平展拓扑中可以唯一地粘合, 命题~\bookref{3.5.13}~的论证可推广, 从而说明 $[U/G]$ 是叠.
最后, 说 $U\to[U/G]$ 是主 $G$-丛, 意味着它可由概形表出, 且沿从概形出发的每个映射 $T\to[U/G]$ 作基变换后都是主 $G$-丛.
\end{remark}
\begin{proof}
我们使用自然投影 $U\to[U/G]$, 它通过 $2$-Yoneda 引理 (\bookref{3.4.21}) 对应于平凡主 $G$-丛 $p_2:G\times U\to U$, 以及由乘法给出的 $G$-等变映射 $\sigma:G\times U\to U$.
为证明 $U\to[U/G]$ 是主 $G$-丛, 设从 $S$-概形出发的态射 $T\to[U/G]$ 由主 $G$-丛 $P\to T$ 及 $G$-等变映射 $P\to U$ 分类. 根据习题~\bookref{3.4.37}, 有 Cartesian 图
\[
\begin{tikzcd}
P\arrow[r]\arrow[d]\arrow[dr,phantom,"\square"]&U\arrow[d]\\
T\arrow[r]&{[U/G]}.
\end{tikzcd}
\]
由于每个基变换都是主 $G$-丛, $U\to[U/G]$ 也是主 $G$-丛. 若 $U$ 是概形, 则映射 $U\to[U/G]$ 是光滑呈示. 一般情形下, 取由概形 $U'$ 给出的平展呈示 $U'\to U$, 则复合 $U'\to U\to[U/G]$ 给出光滑呈示.
\end{proof}
\begin{example}[$B\mathbb G_m$, $B\mathrm{GL}_n$ 与 $B\mathrm{PGL}_n$]\label{src:4.1.12}
由于主 $\mathbb G_m$-丛对应于线丛 (习题~\bookref{B.1.54}), 分类叠 $B\mathbb G_m$ 等价于由二元组 $(S,V)$ 组成的范畴, 其中 $S$ 是概形, $V$ 是 $S$ 上的线丛. 类似地, $B\mathrm{GL}_n$ 是由二元组 $(S,V)$ 组成的叠, 其中 $V$ 是概形 $S$ 上秩为 $n$ 的向量丛 (习题~\bookref{B.1.58}).
分类叠 $B\mathrm{PGL}_n$ 可等价地用 $S$ 上的主 $\mathrm{PGL}_n$-丛, Brauer--Severi 概形或 Azumaya 代数来描述 (见习题~\bookref{B.1.68}~和~\bookref{B.1.70}).
\end{example}
\begin{exercise}[$B\mathrm O(q)$]\label{src:4.1.13}
设 $k$ 是特征 $\operatorname{char}(k)\neq2$ 的域. 对 $n$ 维向量空间 $V$ 上的非退化二次型 $q$, 正交群 $\mathrm O(q)$ 是 $\mathrm{GL}(V)$ 中由保持 $q$ 的矩阵组成的子群. 若 $q,q'$ 是非退化二次型, 证明 $B\mathrm O(q)\cong B\mathrm O(q')$, 尽管 $\mathrm O(q)$ 与 $\mathrm O(q')$ 可能不同构.
\end{exercise}
我们得到一个理想的推论: 代数空间的范畴在对有限群的自由作用取商之下封闭. 回顾例~\bookref{1.5.8}, 有限群在概形上的自由作用的商未必作为概形存在.
\begin{corollary}\label{src:4.1.14}
设 $G$ 是有限抽象群, 将其视为概形 $S$ 上的群概形. 若 $G$ 在 $S$ 上的代数空间 $U$ 上自由作用, 则商层 $U/G$ 是代数空间.
\end{corollary}
\begin{proof}
由于作用自由, 商叠 $[U/G]$ 等价于一个层, 我们将其记为 $U/G$ (见习题~\bookref{3.4.20}). 定理~\bookref{4.1.10}~说明 $U/G$ 是代数叠, 且 $U\to U/G$ 是主 $G$-丛, 因而特别地有限, 平展, 满射且可由概形表出. 取由概形给出的平展呈示 $U'\to U$, 则复合 $U'\to U\to U/G$ 给出 $U/G$ 的平展呈示.
\end{proof}
\begin{remark}\label{src:4.1.15}
此外, 这说明代数空间的范畴本身在对有限群的自由作用取商之下封闭, 因而我们不必进一步扩大所用的范畴.
\end{remark}
\begin{exercise}[简单]\label{src:4.1.16}
设 $G\to S$ 是作用于 $S$-概形 $X,Y$ 的光滑仿射群概形. 证明 $G$-等变态射 $X\to Y$ 诱导代数叠态射 $[X/G]\to[Y/G]$; 反之, $[X/G]\to[Y/G]$ 由 $G$-等变态射诱导当且仅当 $[X/G]\to[Y/G]$ 是 $BG$ 上的态射.
\end{exercise}
